\documentclass[11pt]{article}
\usepackage[margin=1in]{geometry}
\usepackage{amsmath,amssymb,amsthm,booktabs,array}
\usepackage[T1]{fontenc}
\usepackage[colorlinks=true,allcolors=blue]{hyperref}
\newtheorem{theorem}{Theorem}
\newtheorem{proposition}[theorem]{Proposition}
\newtheorem{lemma}[theorem]{Lemma}
\newtheorem{corollary}[theorem]{Corollary}
\newcommand{\R}{\mathbb R}
\newcommand{\Z}{\mathbb Z}
\newcommand{\C}{\mathbb C}
\newcommand{\ind}{\mathbf 1}
\newcommand{\leanfile}{\href{https://github.com/the-omega-institute/trureturing/blob/1ad50310295853e27713693f182e8eac94296739/D5/S3/Fourier/MO511715Source.lean}{\texttt{MO511715Source.lean}}}
\title{A negative answer to mathex's cutoff-series question beyond Ghafari's masks}
\author{Haobo Ma\\\small ChronoAI Pte Ltd, Singapore
\and Wenlin Zhang\\\small National University of Singapore\\\small ChronoAI Pte Ltd, Singapore}
\date{19 November 2026}
\begin{document}
\maketitle
\begin{abstract}
For mathex's question \emph{Bound of series}, MathOverflow 511715,
revision 2 (26 May 2026), final displayed inequality, we give a negative answer to the requested
uniform cutoff estimate in the Euclidean interpretation, complementing
Ghafari's bound for a related Littlewood--Paley-mask series
(2026, Appendix H, Lemma H.0.1): one admissible
mask isolates a nonzero output, and disjoint lattice blocks force
logarithmic divergence as the cutoff scale tends to zero.
For each nonzero output, the same symmetric mask works at every positive
time and for every fixed real or complex compactly supported cutoff
continuous at zero and equal to one there. No finite constant and fixed real
exponent can satisfy the proposed estimate even on any bounded interval
of positive scales reaching zero.
For smooth cutoffs, a corrected Lipschitz estimate holds uniformly in the
mask and output when time is bounded and the cutoff scales stay away from zero;
zero modes show that the time restriction cannot simply be dropped.
\end{abstract}

\section{The question and its precise interpretation}

Mathex's question~\cite{mathex} asks ``How to establish'' a uniform
H\"older estimate after subtracting the zero-output value of a
five-dimensional cutoff series.
The support condition on the mask allows it to isolate a single nonzero
output, so that subtraction need not cancel anything. We retain the full
time integral and all lattice modes and show that this freedom causes
divergence even at fixed positive time.

The closest positive result concerns a prescribed Littlewood--Paley mask.
Ghafari~\cite{ghafari}, Appendix H, defines this mask on printed p.~93
and the same series on p.~108 (with renamed summation indices).
Lemma H.0.1, pp.~108--111, bounds its output difference, uniformly in
$\epsilon>0$ and bounded time, by
$(1+|n|^4)^{\gamma/2}$ for every $\gamma>0$; the proof yields
$1+\log(2+|n|)$. It uses a smooth frequency extension with
$|\nabla_p\psi(p,s)|\lesssim(1+|s|)^{-1}$ on comparable frequencies
(p.~109). This is a bound on the output difference at one scale, not
the cutoff-increment estimate uniform in $n$ asked here.
The one-way support condition alone supplies neither that derivative
control nor its cancellation. Corollary~\ref{cor:regularity} makes
this distinction explicit, so the obstruction does not contradict
Ghafari's result.

All vector norms below are Euclidean. Let $F=\R$ or $\C$; scalar
absolute values mean absolute value or complex modulus. Fix
$\zeta:\R^5\to F$, compactly supported with $\zeta(0)=1$.
The negative result requires continuity only at zero; smoothness is
needed for the corrected estimate in Section~\ref{sec:correction}.
Write $D=2\pi$. A mask is a function
$\psi:(\Z^5)^2\to[0,1]$ satisfying
\begin{equation}\label{eq:mask}
 \psi(k,l)\ne0\quad\Longrightarrow\quad
 \frac9{64}|k|\le |l|\le\frac{64}9|k|.
\end{equation}
For $q\ge0$ and $L,A,B\ge0$, define
\begin{align}
 I(q;L,A,B)&=\int_0^q e^{-L(q-\alpha)}
 \int_0^\alpha\int_0^\alpha
 e^{-A(q+\alpha-2v_1)-B(q+\alpha-2v_2)}\,dv_2\,dv_1\,d\alpha,
 \label{eq:integral}\\
 K_\psi(q,n,m_1,m_2)&=D^2|n-m_1|^2\psi(n-m_1,m_1)
 \notag\\[-4pt]
 &\hspace{12mm}\cdot I\bigl(q;D^4|n-m_1|^4,D^4|m_1-m_2|^4,D^4|m_2|^4\bigr).
 \label{eq:kernel}
\end{align}
The iterated inner integral is the integral over the full square
$[0,\alpha]^2$ with product Lebesgue measure. For $\epsilon>0$, put
\begin{align}
 \phi_\epsilon(q,n)&=\sum_{m_1,m_2\in\Z^5}
 K_\psi(q,n,m_1,m_2)
 |\zeta(\epsilon(m_1-m_2))|^2|\zeta(\epsilon m_2)|^2,
 \label{eq:phi}\\
 \Delta_\epsilon(r,u)&=\phi_\epsilon(r,u)-\phi_\epsilon(r,0).
 \label{eq:delta}
\end{align}
Every sum in~\eqref{eq:phi} has finite support: choose $R>0$ with
$\zeta(x)=0$ for $|x|>R$. A nonzero summand forces both
$|m_1-m_2|\le R/\epsilon$ and $|m_2|\le R/\epsilon$.
The corresponding integer vectors form finite sets; the map
$(a,b)\mapsto(m_1,m_2)=(a+b,b)$ is a bijection.

The proposed estimate is
\begin{equation}\label{eq:asked}
 |\Delta_{\epsilon_2}(r,u)-\Delta_{\epsilon_1}(r,u)|
 \le C|\epsilon_2-\epsilon_1|^\delta,
\end{equation}
uniformly in $r,u,\epsilon_1,\epsilon_2,\psi$. We allow the constants
to depend on the individually fixed cutoff. The question does not specify
the vector norm or scalar field, or the signs of $C,\delta$.
The conventional reading uses $C\ge0$, $\delta>0$, and $r\ge0$;
our negative result covers all finite real $C,\delta$ under the
Euclidean real and complex readings stated above.

\section{A fixed-output obstruction}

For each $v\in\Z^5\setminus\{0\}$ choose the single mask
\begin{equation}\label{eq:selected}
 \psi_v(k,l)=\ind_{\{k+l=v\}}
 \ind_{\{(9/64)|k|\le |l|\le(64/9)|k|\}}.
\end{equation}
This mask satisfies~\eqref{eq:mask} and is independent of the cutoff,
time, and scale. At $(n-m_1,m_1)$ its first indicator forces $n=v$.
Consequently, for this mask,
\begin{equation}\label{eq:isolates}
 \phi_\epsilon(q,n)=0\quad(n\ne v),\qquad
 \Delta_\epsilon(q,v)=\phi_\epsilon(q,v).
\end{equation}
Define the positive constants
\begin{equation}\label{eq:constants}
 \eta=\frac1{20000D^4},\qquad
 \kappa=D^2e^{-1}\eta^3,\qquad c=\frac\kappa{16}.
\end{equation}

\begin{theorem}\label{thm:main}
For every compactly supported $\zeta:\R^5\to F$ continuous at zero
with $\zeta(0)=1$, there is $\rho\in(0,1]$ with the following property.
Fix any $q>0$ and $v\in\Z^5\setminus\{0\}$, and choose a nonnegative
integer $j_0$ such that
\[
 2^{j_0}\ge\max\{1,|v|\},\qquad \eta\,2^{-4j_0}\le q/2.
\]
For the mask $\psi_v$ and every $0<\epsilon\le\rho/(5\cdot2^{j_0})$,
\begin{equation}\label{eq:loggrowth}
 \Delta_\epsilon(q,v)\ge
 c\left(1+\left\lfloor\log_2\frac{\rho}{5\epsilon}\right\rfloor-j_0\right).
\end{equation}
In particular $\Delta_\epsilon(q,v)\to+\infty$ as $\epsilon\downarrow0$.
For $q=1$ and $v=e=(1,0,0,0,0)$ one may take $j_0=0$.
\end{theorem}

\begin{lemma}\label{lem:time}
If $q>0$, $0<h\le q/2$, and $(L+3A+3B)h\le1$, then
$I(q;L,A,B)\ge e^{-1}h^3$.
\end{lemma}
\begin{proof}
Restrict the integral to $q-h\le\alpha\le q$ and
$\alpha-h\le v_1,v_2\le\alpha$. Since $h\le q/2$, this region lies
in the original integration domain. On it,
$0\le q-\alpha\le h$ and $0\le q+\alpha-2v_i\le3h$.
The full integrand is at least $e^{-1}$ and the region has volume
$h^3$. Nonnegativity permits the restriction.
\end{proof}

\begin{proof}[Proof of Theorem~\ref{thm:main}]
For each integer $N\ge\max\{1,|v|\}$, let
$B_N=[N,2N)^5\cap\Z^5$. It has $N^5$ elements. Given
$a,b\in B_N$, put $m=a+b$ and $k=v-m$.
Coordinatewise bounds and the triangle inequality give
\[
 |a|,|b|\le\sqrt{20}N\le5N,\qquad
 4N\le|m|\le\sqrt{80}N\le9N,\qquad
 3N\le4N-|v|\le|k|\le9N+|v|\le10N.
\]
In particular
$(9/64)|k|\le(90/64)N<4N\le|m|$, whereas
$|m|\le9N<(64/9)3N\le(64/9)|k|$.
Thus $\psi_v(k,m)=1$ throughout this block.

For $N\ge2^{j_0}$ take $h=\eta/N^4$. Then $0<h\le q/2$, and for the three rates
$L=D^4|k|^4$, $A=D^4|a|^4$, $B=D^4|b|^4$,
\[
 (L+3A+3B)h
 \le\frac{10^4+6\cdot5^4}{20000}=\frac{11}{16}<1.
\]
Lemma~\ref{lem:time}, the actual squared-frequency prefactor,
and $|k|\ge N$ therefore imply
\begin{equation}\label{eq:blockkernel}
 K_{\psi_v}(q,v,a+b,b)
 \ge D^2N^2 e^{-1}\frac{\eta^3}{N^{12}}
 =\frac\kappa{N^{10}}.
\end{equation}

Continuity at zero supplies $0<\rho\le1$ such that
$|\zeta(x)|\ge1/2$ whenever $|x|\le\rho$.
For $0<\epsilon\le\rho/(5\cdot2^{j_0})$, set
$J=\lfloor\log_2(\rho/(5\epsilon))\rfloor\ge j_0$.
For every $j_0\le j\le J$ and $a,b\in B_{2^j}$,
$|\epsilon a|,|\epsilon b|\le5\epsilon2^j\le\rho$.
The product of the two squared cutoff factors is at least $1/16$.
There are $2^{10j}$ ordered pairs in this block, so
\eqref{eq:blockkernel} gives mass at least $\kappa/16=c$.

The boxes $B_{2^j}$ are pairwise disjoint: their first-coordinate
intervals are disjoint half-open dyadic intervals. Hence their pair
blocks are disjoint, and the bijection $(a,b)\mapsto(a+b,b)$ preserves
disjointness. All other summands are nonnegative. Summing these $J-j_0+1$
blocks and using~\eqref{eq:isolates} proves~\eqref{eq:loggrowth}.
In particular, at $\epsilon_J=\rho/(5\cdot2^J)$,
\begin{equation}\label{eq:dyadic}
 \Delta_{\epsilon_J}(q,v)\ge c(J-j_0+1)\qquad(J\ge j_0).
\end{equation}
\end{proof}

\begin{corollary}\label{cor:interval}
For a cutoff as in Theorem~\ref{thm:main}, fix $q>0$, $v\ne0$,
and $E>0$. No finite real pair $C,\delta$
satisfies~\eqref{eq:asked} for all distinct
$\epsilon_1,\epsilon_2\in(0,E]$, even for the single symmetric mask $\psi_v$.
\end{corollary}
\begin{proof}
The support inequalities in~\eqref{eq:selected} are equivalent after
interchanging $k,l$, since $9/64$ and $64/9$ are reciprocal; hence
$\psi_v(k,l)=\psi_v(l,k)$. Let $M=\phi_E(q,v)$ and
$H=\max\{(E/2)^\delta,E^\delta\}$, both finite.
Since $\epsilon_J\downarrow0$, the gap $|E-\epsilon_J|$ lies in
$[E/2,E]$ for large $J$. Thus
\[
 C|E-\epsilon_J|^\delta\le\max\{C,0\}H
 \quad\text{for every real }C,\delta.
\]
Choose $J$ so large that
$c(J-j_0+1)>M+\max\{C,0\}H$. Equations~\eqref{eq:isolates}
and~\eqref{eq:dyadic} yield
\[
 |\Delta_E(q,v)-\Delta_{\epsilon_J}(q,v)|
 \ge\phi_{\epsilon_J}(q,v)-M
 >C|E-\epsilon_J|^\delta.
\]
Both scales use the same cutoff, time, output, and mask.
\end{proof}

\begin{corollary}\label{cor:regularity}
Fix $v\ne0$. No interpolation $\Psi:\R^5\times\Z^5\to\R$
of $\psi_v$, continuously differentiable in its first argument,
has a finite constant $A$ such that
\[
 |\nabla_p\Psi(-s+\theta v,s)|\le \frac{A}{1+|s|}
 \quad\text{for all }|s|>2|v|,\ 0\le\theta\le1.
\]
\end{corollary}
\begin{proof}
Take $s=(2N,2N,2N,2N,2N)$ for integers $N\ge\max\{1,|v|\}$.
The block bounds in the proof of Theorem~\ref{thm:main}, with
$a=b=(N,N,N,N,N)$, give $\psi_v(v-s,s)=1$;
also $\psi_v(-s,s)=0$ because $v\ne0$.
The mean value inequality along the segment would imply
$1\le A|v|/(1+|s|)$ for arbitrarily large $N$, a contradiction.
\end{proof}

\section{A corrected estimate away from zero scale}\label{sec:correction}

The obstruction concerns scales tending to zero. On a bounded time
interval the following uniform estimate survives.

\begin{proposition}\label{prop:correction}
Let $\zeta\in C_c^\infty(\R^5;F)$, $\sigma>0$, and $0\le T<\infty$.
There exists $L_{\zeta,\sigma,T}<\infty$ such that, for every mask
satisfying~\eqref{eq:mask}, every $u\in\Z^5$, $0\le r\le T$,
and $\epsilon_1,\epsilon_2\ge\sigma$,
\[
 |\Delta_{\epsilon_2}(r,u)-\Delta_{\epsilon_1}(r,u)|
 \le L_{\zeta,\sigma,T}|\epsilon_2-\epsilon_1|.
\]
\end{proposition}
\begin{proof}
Choose $R>0$ with $\zeta(x)=0$ for $|x|>R$, and let
$B_0=R/\sigma$, $S=\{a\in\Z^5:|a|\le B_0\}$.
For every $\epsilon\ge\sigma$, only $a,b\in S$ occur in the
reindexed sum. Put
$M_0=\sup_x|\zeta(x)|$ and $M_1=\sup_x\|d\zeta_x\|$,
where the derivative is real-linear and its norm is the operator norm.
Both are finite by smooth compact support.
For fixed $a,b\in S$, the function
$w_{a,b}(\epsilon)=|\zeta(\epsilon a)|^2|\zeta(\epsilon b)|^2$
satisfies
\[
 |w'_{a,b}(\epsilon)|
 \le2M_0^3M_1(|a|+|b|)\le4M_0^3M_1B_0=:W_*.
\]
If the mask is nonzero,~\eqref{eq:mask} implies
$|n-a-b|\le(64/9)|a+b|\le(128/9)B_0$.
Every exponential in~\eqref{eq:integral} is at most one on its
integration domain, so $I(q;L,A,B)\le q^3/3$.
Therefore, for all $q\in[0,T]$ and all $n\in\Z^5$,
\[
 0\le K_\psi(q,n,a+b,b)
 \le\left(D\frac{128}{9}B_0\right)^2\frac{T^3}{3}=:K_*;
\]
if the mask vanishes this inequality is immediate.
The mean value inequality, applied to each of the $|S|^2$ terms,
gives a Lipschitz constant $|S|^2K_*W_*$ for
$\phi_\epsilon(q,n)$, uniformly in $q,n,\psi$.
The difference of the two outputs has constant
$L_{\zeta,\sigma,T}=2|S|^2K_*W_*$, as required.
This also includes $T=0$, when all terms vanish.
\end{proof}

\begin{corollary}\label{cor:time}
For each cutoff in Theorem~\ref{thm:main}, there are two fixed positive
scales $\epsilon_a<\epsilon_b$ and $d>0$ such that, with
$\psi_e$ and $e=(1,0,0,0,0)$,
\[
 |\Delta_{\epsilon_b}(q,e)-\Delta_{\epsilon_a}(q,e)|\ge dq
 \qquad(q\ge2).
\]
Thus a positive lower bound on the scales alone does not give a
time-uniform replacement for~\eqref{eq:asked}.
\end{corollary}
\begin{proof}
Choose $R>0$ with $\zeta(x)=0$ for $|x|>R$, and take
$\epsilon_a=\rho/4$, $\epsilon_b=R+1$.
At scale $\epsilon_b$ only $a=b=0$ can contribute; then the
mask or the frequency prefactor vanishes, so $\phi_{\epsilon_b}(q,n)=0$.
At scale $\epsilon_a$ retain the one term $a=0$, $b=2e$.
Its output frequency is $k=-e$; the mask equals one, and its cutoff
weight is $|\zeta(0)|^2|\zeta(2\epsilon_a e)|^2\ge1/4$.
For $L=D^4$ and $B=16D^4$, direct integration of the two inner
variables gives
\[
 I(q;L,0,B)=\int_0^q e^{-(L+B)(q-\alpha)}
             \frac{\alpha(1-e^{-2B\alpha})}{2B}\,d\alpha.
\]
For $q\ge2$, restrict to $q-1\le\alpha\le q$. There
$\alpha\ge q/2$, $\alpha\ge1$, and the integrand is at least
$q e^{-(L+B)}(1-e^{-2B})/(4B)$.
Multiplying by the prefactor $D^2$ and cutoff weight $1/4$ proves
the claim with $d=D^2e^{-(L+B)}(1-e^{-2B})/(16B)>0$.
All other terms are nonnegative, and~\eqref{eq:isolates} removes
the zero-output value. These two scales lie in $[\epsilon_a,\infty)$;
their separation is fixed while $dq$ is unbounded.
\end{proof}

\section{Coverage and formal statements}

The source has one question, no linked mathematical source, and one
comment asking for motivation rather than posing another problem.
The consequences below that are not in that question are marked as such.

\begin{center}
\begin{tabular}{@{}>{\raggedright\arraybackslash}p{0.37\linewidth}>{\raggedright\arraybackslash}p{0.55\linewidth}@{}}
\toprule
Statement & Conclusion\\
\midrule
Posed estimate~\eqref{eq:asked} & Answered: no under the stated
Euclidean conventions; corrected by Proposition~\ref{prop:correction}.\\
Every fixed positive time and nonzero output & Not posed; logarithmic
obstruction proved in Theorem~\ref{thm:main}.\\
Symmetric masks; scales in $(0,E]$ & Not posed; obstruction proved in
Corollary~\ref{cor:interval}.\\
Positive lower scale, bounded time & Not posed; corrected Lipschitz
estimate proved in Proposition~\ref{prop:correction}.\\
Positive lower scale, unbounded time & Not posed; linear obstruction
proved in Corollary~\ref{cor:time}.\\
Sole motivation comment & Explanatory, not a mathematical assertion.\\
Ghafari, Lemma H.0.1 & Prior proof cited: structured-mask output bound;
not the source's cutoff-increment question.\\
Frequency derivative control for $\psi_v$ & Not posed; obstruction
proved in Corollary~\ref{cor:regularity}.\\
Additional cancellation or partition identities for the mask & Not covered;
such identities would specify a different class from~\eqref{eq:mask}.\\
\bottomrule
\end{tabular}
\end{center}

\begin{sloppypar}
In the pinned Lean source~\cite{lean}, within namespace
\texttt{D5.S3.Fourier.MO511715Source}, the declarations
\texttt{cutoff\_finite\_support}, \texttt{time\_integral\_ge},
and \texttt{kernel\_block\_lower} give the finite-support,
time-subregion, and block-kernel statements. The latter two are at
$q=1$, and the block statement uses output $e$ and mask $\psi_e$.
The declaration \texttt{phi\_dyadic\_growth} gives~\eqref{eq:dyadic}
at $q=1$, $v=e$, $j_0=0$ for both scalar fields.
The declaration \texttt{result} proves the negation of
\texttt{claim}, the conjunction of the real and complex universal
smooth-cutoff estimates. Its proof uses a real smooth bump;
the stronger fixed-cutoff, all-real-exponent failure is the intermediate
private lemma \texttt{source\_estimate\_failure}, at time one,
output $e$, with the larger scale fixed at two.
The all-scale floor bound~\eqref{eq:loggrowth}, its extension to arbitrary
positive times and nonzero outputs, Proposition~\ref{prop:correction},
and Corollaries~\ref{cor:interval}, \ref{cor:regularity}, and~\ref{cor:time}
have the written proofs above;
their full statements are not the conclusions of those Lean declarations.
The accompanying \texttt{verify.py} checks finite lattice instances and
exact rational identities used in the block estimates.
\end{sloppypar}

\begin{thebibliography}{3}
\raggedright
\bibitem{mathex}
mathex (MathOverflow user 138491),
\emph{Bound of series}, MathOverflow question 511715,
revision 2, 26 May 2026,
\href{https://mathoverflow.net/revisions/511715/2}{revision text};
\href{https://mathoverflow.net/questions/511715/bound-of-series}{question}.
Source text licensed CC BY-SA 4.0.

\bibitem{ghafari}
Joe Ghafari, \emph{Paracontrolled approach to the stochastic Cahn--Hilliard
equation}, Ph.D. thesis, University of Ottawa, 2026,
\href{https://arxiv.org/abs/2609.02197v1}{arXiv:2609.02197v1},
2 September 2026. Appendix H, pp.~93 and 108--111, Lemma H.0.1.

\bibitem{lean}
\emph{Five-dimensional cutoff series}, Lean source
\leanfile, commit
\texttt{1ad50310295853e27713693f182e8eac94296739}.
\end{thebibliography}

\smallskip
\noindent\small This note is licensed under
\href{https://creativecommons.org/licenses/by/4.0/}{Creative Commons Attribution 4.0
International (CC BY 4.0)}.
\end{document}
